%! Problems with Sampling — Unit 3, Lesson 3.3 — Warm-Up
% ONE PAGE, blank and key. Three items; do not add a fourth.
% GRADUAL RELEASE: the warm-up activates prior knowledge before the guided notes.
%   Item 1 — spiral to Lesson 3.2: name two random sampling methods (SRS, stratified).
%   Item 2 — Maya's lunch-table survey: did chance decide, and who is left out? Seeds row 1
%            (convenience) and row 2 (who can never be picked).
%   Item 3 — two wordings of one question: which gets more "yes"? Seeds row 3 (response bias).
% Every display is given; nothing is sketched.
\documentclass[12pt]{article}
\usepackage{apstats-article}
\usepackage{apstats-boxes}

\begin{document}

\pageheader{Unit 3, Lesson 3.3}{Warm-Up}
% NAMESTRIP: no \namedateperiod here — the name row lives on the cover only.

\vspace{0.28in}

\begin{enumerate}[leftmargin=1.9em, itemsep=0.42in]

  \item \textbf{From Lesson 3.2.} A principal wants to survey 60 of her school's 1{,}200
        students. Name each sampling method.

        \vspace{12pt}
        \textbf{(a)} She numbers all 1{,}200 students and has a computer pick 60 at random.

        \vspace{8pt}
        Method: \blank{6.0cm}

        \vspace{12pt}
        \textbf{(b)} She picks 15 students at random from each of the four grades.

        \vspace{8pt}
        Method: \blank{6.0cm}

  \item Maya wants to know how many hours a week the students at her school work at paid jobs.
        She asks the 12 friends at her lunch table. Did chance decide who was asked? Which
        students might her sample leave out?

        \vspace{12pt}
        \writeline\\[18pt]\writeline

  \item Which version of this question would get more ``yes'' answers? Write A or B, then
        explain.

        \vspace{10pt}
        \begin{center}
        \renewcommand{\arraystretch}{1.4}
        \begin{tabular}{@{} c p{0.72\linewidth} @{}}
        \textbf{(A)} & ``Should our school ban phones in class?'' \\
        \textbf{(B)} & ``Should our school ban phones in class so that students can focus and
                       learn?'' \\
        \end{tabular}
        \end{center}
        \vspace{6pt}
        \writeline\\[18pt]\writeline\\[18pt]\writeline

\end{enumerate}

\end{document}
